\label{sec:limitations}
\textit{Data.} The DOT link set is sparse and highway-heavy; the full-year link enumeration timed out server-side, so
the study uses the 26 Manhattan links active since 2026-09-20, of which seven are dead feeds, three more have under
3\,\% coverage, three lose their feed before or during the test partition (NOT\_EVALUATED) and 16 are modelled; about
half of all raw records are zero-speed no-measurement rows, treated as missing. Monitored segments do not tile the
highways and the published polylines stop short of the named link ends, so the graph needs name-based continuity
rules and remains sparse (16 directed edges, 6 isolated nodes before eligibility). The eligibility threshold was
lowered from the planned 0.6 to 0.5 training coverage after the coverage audit and before any result on the 0.6 set
existed; both are recorded. Operational weather is Open-Meteo's historical-forecast product, a re-forecast proxy
rather than an archived agency forecast. Collisions lag publication by months and are retrospective only.

\textit{Modelling and evaluation.} All deep models were trained under a recorded budget of 4{,}096 windows per epoch
with early stopping, one seed and one chronological split; the only variance estimate is a day-block bootstrap on
paired differences, and the proposed model's margin over the GRU and temporal Transformer is inside what seeds could
change. Split conformal intervals carry a marginal guarantee under exchangeability, which traffic series violate; the
measured coverage under congestion is 0.76--0.80. Explanations describe the trained model, with a background drawn
from the training partition; neighbour sensitivity perturbs the graph in ways the model never saw during training,
and only incoming edges can change a link's own prediction.

\textit{Knowledge, language model and agents.} The approved corpus is small and is not NYC DOT's signal-timing
policy; the project-defined policy is labelled as such. 171 of 1{,}633 chunks exceed the encoder's 256-token window.
The retrieval labels were drafted by the project without external review, so recall and MRR are an internal
consistency check. The language model is a 7B quantised model served locally; its advisories are only as good as the
evidence and the constraints that follow it. The advisory evaluation is functional (seven cases, four policies, 19
runs), not statistical, and on the mapped corridor the model proposed no signal change, so the simulation acceptance
path was not exercised by a real advisory. Two defects were found during the run and fixed before the reported
numbers were produced (the bootstrap day index, and a missing field that had silently prevented every LLM call in a
first agent pass); both superseded runs are kept as local evidence and described in \texttt{progress.md}.

\textit{Simulation.} The SUMO testbed is a scenario built from OSM geometry with declared demand and default signal
programs, not a calibrated digital twin, and contains no pedestrian demand; the mapped DOT links are not the
forecastable ones, so the controller comparison is a transfer demonstration. \textit{Engineering.} Data download,
training, LLM inference and simulation ran on the authors' workstation (torch 2.11.0+cu128, a recorded deviation from
the lock file); the development container could not reach the data portals or the model server, and the container
images were built and smoke-tested on the workstation only. Model weights, prediction arrays, raw data and the FAISS
index are not in the repository; their hashes and run ids are. The complete, maintained list is in
\texttt{reports/limitations.md}.
